\worldchapter{Charaktererschaffung}{create}
\label{ch:create}

\begin{multicols}{2}

Zur Erstellung eines Charakters werden Charakterschablonen ausgewählt, die dessen Lebensstationen, Talente und Interessen abbilden (siehe \cref{ch:appendix-templates}). Jede Schablone kann die Attribute und Fertigkeiten des Charakters verändern und Wissen sowie Schatten mit sich bringen.

\wbif{tirakan}{
Darüber hinaus können Charakterschablonen Spielmechaniken freischalten, wie etwa das Ausführen von Priesteraktionen oder das Erlernen von Zaubern eines Magieursprungs.
}{}

\section{Reputation}

Die Reputation verdeutlicht die Bekanntheit des Charakters. Jedes Abenteuer bringt dem Charakter eine gewisse Menge an Reputation ein. Neue Charaktere starten ihre Laufbahn in der Regel mit 20 Punkten Reputation. Dies kann jedoch von der Spielleiterin festgelegt werden.

Reputation wird verwendet, um Charakterschablonen hinzuzufügen. Jede Schablone kostet eine bestimmte Anzahl an Reputationspunkten.

Charakterschablonen können negative Reputationskosten aufweisen. In diesem Fall erhält der Spieler die Punkte, wenn er die Schablone auswählt. Dies ist beispielsweise bei der Schablone \textit{Säufer} der Fall.

\wbif{phasesix}{
\section{Epochen}

Vor Beginn der Kampagne oder des Abenteuers legt die Spielleiterin fest, in welcher Epoche und mit welchen Erweiterungen gespielt wird. Diese Auswahl legt fest, welche Charakterschablonen, Waffen, Rüstungen und Gegenstände verwendet werden können und ob Magie, Körpermodifikationen oder Priesteraktionen möglich sind.

Mögliche Epochen sind:

\begin{itemize}
    \item Die klassische Antike \faIcon{eye}
    \item Mittelalter, Wikinger und Kreuzzüge \faIcon{chess-rook}
    \item Viktorianisches Zeitalter und der Wilde Westen \faIcon{umbrella}
    \item Imperialismus und Weltkriege \faIcon{flag}
    \item Der kalte Krieg und die 80er \faIcon{plane}
    \item Moderne Zeit \faIcon{microchip}
%    \item Nahe Zukunft
    \item Science Fiction \faIcon{space-shuttle}
\end{itemize}

Optionale Erweiterungen sind:

\begin{itemize}
    \item Magie \faIcon{magic}
    \item Horror \faIcon{spider}
    \item Pantheon \faIcon{ankh}
    \item Körpermodifikationen \faIcon{syringe}
%    \item Fahrzeuge und Dronen \faIcon{car}
\end{itemize}

Bei allen Schablonen, Gegenständen und Waffen sind die Epochen und Erweiterungen aufgelistet, in denen sie gültig sind. Bei den Charakterschablonen, Gegenständen etc. im Anhang werden dafür die Symbole verwendet, die in der Auflistung oben stehen.
}{}


\section{Schablonen wählen}

Eine Charakterschablone stellt eine bestimmte Lebensstation eines Charakters dar. Jede Schablone ist einem der folgenden Bereiche zugeordnet: Bildung, Beruf, Talent, Interessen, Charakter oder Lebensumstände.

Jede Schablone verändert eine kleine Anzahl von Eigenschaften und Fertigkeiten des Charakters positiv oder negativ und kann Wissen oder Schatten mit sich bringen. Darüber hinaus können Schablonen eigene Regeln enthalten, die der Charakter dann übernimmt.

\begin{tcolorbox}
\subsection*{Gelehrter}
\tcblower
\begin{tabular}{@{}ll@{}}
    \textbf{Reputation:} & 10\\
\end{tabular}

\begin{tabular}{@{}ll@{}}
    \textbf{Bildung} & +4\\
    \textbf{Naturkunde} & +1\\
    \textbf{Historie} & +2\\
    \textbf{Kommunikation} & +1\\
\end{tabular}
\end{tcolorbox}

Jede Schablone ist eine bestimmte Anzahl an Reputationspunkten wert. Dies ist die Anzahl der Reputation, die aufgewendet werden muss, um die Schablone in den eigenen Werdegang zu übernehmen.

Die Liste der Schablonen ist in \cref{ch:appendix-templates} zu finden.

\subsection{Basiswerte}

Alle Attribute, Fertigkeiten und sonstigen Werte des Charakters beginnen mit einem einheitlichen Wert. Auf diesen Wert werden die Angaben der Charakterschablonen addiert.

\begin{itemize}
    \item Aktionen: 2
    \item Mindestwurf: 5+
    \item Bonus-, Schicksals- und Wiederholungswürfe: 0
    \item Attribute: 1
    \item Fertigkeiten: 0
    \item Angeborener Schutz: 0
    \item Maximale Wunden: 10
    \wbif{nexus}{}{
    \item Arkana: 0
    \item Zauberpunkte: 0
    }
    \wbif{tirakan}{}{
    \item Maximaler Stress: 10
    \item Basisstress: 0
    \item Biolast: 0
    }
\end{itemize}

\subsection{Abstammung}

Zunächst wählst du bei der Charaktererschaffung die Abstammung deines Charakters. Diese Herkunft beschreibt nicht nur die Kultur, aus der dein Charakter stammt, sie bringt auch die Abstammungsschablone mit sich, die dem Charakter die üblichen Stärken eines Vertreters deines Volkes gibt.

Es kann nur eine Abstammungsschablone gewählt werden, diese kostet keine Reputation.

Die verfügbaren Abstammungsschablonen sind im \cref{ch:appendix-templates} aufgelistet.

Die Abstammungsschablone wird beim Werdegang notiert, und die angegebenen Modifikationen werden auf die Werte des Charakters addiert.

\subsection{Weitere Schablonen}

Nun können nach Belieben weitere Schablonen ausgewählt werden, bis keine Reputation mehr verfügbar ist. Es können Schablonen aus allen Kategorien beliebig zusammengestellt werden. So kann der Charakter einen oder mehrere Berufe haben oder auch gar keinen.

Die angegebenen Modifikationen der jeweiligen Schablone werden auf die Werte des Charakters addiert. Zusätzlich werden Wissen, Schatten und weitere Regeln der Schablone dem Charakterbogen hinzugefügt.

Dabei können alle Werte auch negativ werden (siehe \cref{ch:rolls}).

\subsection{Übrige Reputation}

Ist der Spieler mit der Zusammenstellung der Schablonen zufrieden, kann er den Charakter einfach als \textit{fertig} deklarieren. Sollte noch nicht ausgegebene Reputation übrig sein, so werden diese der \textit{Reputation} des Charakters hinzugefügt (siehe \cref{ch:advancement}). Es geht also keine Reputation verloren.

\section{Kontakte und Sprachen}

Nachdem die Charakterschablonen festgelegt sind, können Sprachen und Kontakte des Charakters festgelegt werden.

\subsection{Kontakte}

Kontakte sind Bekanntschaften oder Verbindungen, die der Charakter bereits vor der Kampagne hatte. Kontakte werden mit Namen und Beschreibung festgehalten, und können nach Belieben ausgedacht werden.

Ein neuer Charakter kann eine bestimmte Anzahl an Kontakten haben, die sich nach der Summe der Attribute \textit{Charme} und \textit{Attraktivität} richtet.

Die Kontakte werden auf dem Charakterbogen festgehalten.

\subsection{Sprachen}

Ein neuer Charakter kann eine bestimmte Anzahl von Sprachen entsprechend der Summe seiner Attribute \textit{Bildung} und \textit{Logik} erlernen. Dies können beliebige Sprachen sein. Beträgt die Summe der Attribute 0 oder weniger, beherrscht der Charakter selbst seine Muttersprache nur eingeschränkt.

Die Sprachen werden auf dem Charakterbogen festgehalten.

\section{Ausrüstung}

Nachdem die Charakterwerte durch die Schablonen festgelegt wurden, kann der Charakter noch mit Ausrüstung ausgestattet werden. Die Spielleiterin legt für die Kampagne oder das Abenteuer ein Startkapital für die Charaktere fest.

\wbif{tirakan}{
Das Startkapital beträgt üblicherweise 2.000 Gulden.
}{
Das Startkapital beträgt üblicherweise 2.000 Einheiten der üblichen Währungseinheit, beispielsweise Euro.
}

Dieses Startkapital kann dafür verwendet werden, um Ausrüstung, wie etwa Waffen, Rüstungen und Gegenstände zu kaufen. Näheres dazu findet sich im Kapitel \cref{ch:equipment}.

\subsection{Ausrüstungsgegenstände}

\cref{ch:appendix-weapons}, \cref{ch:appendix-riotgear} und \cref{ch:appendix-items} können nun für das Startkapital erworben werden. Beliebige Ausrüstungsgegenstände können mit ihren Werten auf dem Charakterbogen notiert, und der Preis vom Startkapital abgezogen werden.

\subsection{Vermögen}

Alles Startkapital, das nicht für Waffen, Rüstungen und Ähnliches ausgegeben wurde, wird zum Vermögen des Charakters.

\wbif{nexus}{}{
\section{Zauber}

\wbif{tirakan}{}{
Wird im Abenteuer oder der Kampagne die Magie-Erweiterung verwendet, so kann der Charakter auch Zauber erlernen.
}

Charakterschablonen bieten \textit{Zauberpunkte} und ermöglichen das Erlernen von Zaubern eines bestimmten \textit{Ursprungs}.

Hat der Charakter nun durch die Wahl von Charakterschablonen beides erhalten, so kann er für die Zauberpunkte Zauber wählen, die er beherrscht.

Zauber werden ähnlich wie Schablonen für Punkte erworben. Hierfür werden Zauberpunkte verwendet. Jeder Zauber hat bestimmte Kosten, für die er dem Charakterbogen hinzugefügt werden kann (siehe \cref{ch:appendix-spells}). Es können nur Zauber der Ursprünge gewählt werden, die der Charakter durch Charakterschablonen freigeschaltet hat. Näheres dazu findet sich im Kapitel \cref{cha:chapter-magic}.
}

\wbif{tirakan}{}{
\section{Körpermodifikationen}

\wbif{nexus}{
Im Gegensatz zur normalen Bevölkerung hat die NEXUS-Organisation Zugang zu teilweise außerirdischen Körpermodifikationen. Dabei handelt es sich um mechanische oder elektronische Elemente, die ein Charakter in seinen Körper integrieren kann.
}{
Wird mit der Erweiterung \textit{Körpermodifikationen} gespielt, so können für das Startkapital \cref{ch:appendix-bodymodifications} erworben und eingebaut werden.
}

Köperermodifikationen können zu Beginn für das Startkapital erworben werden. Hierbei sind die Regeln (siehe \cref{ch:bodymodifications}) für Körpermodifikationen zu berücksichtigen, so muss beispielsweise genügend Energie für die Verbraucher verfügbar sein.

Der Vorgang der Integration durch einen Arzt entfällt bei der Charaktererschaffung, die Körpermodifikationen können einfach auf dem Charakterbogen vermerkt werden.
}
\end{multicols}
